% sec:ad-reverse — Reverse-Mode AD
\section{Reverse-Mode AD\starmark{}}
\label{sec:ad-reverse}

Forward-mode AD propagates derivatives alongside values in a
single forward pass through the computation.  Each seeded
evaluation produces one column of the Jacobian --- one
directional derivative --- so computing the full Jacobian of
a function $f : \mathbb{R}^n \to \mathbb{R}^m$ requires $n$
passes.  For the metric function ($n = 4$, $m = 10$), four
passes suffice and the cost is modest.  But gradient-based
optimisation of spacetime parameters --- fitting a black hole
mass and spin to observational data, for example --- inverts
this ratio: the objective is a single scalar ($m = 1$), and
the parameters may number in the dozens or hundreds ($n \gg 1$).
Forward mode would require one pass per parameter.
Reverse-mode AD solves this by propagating derivatives
backward through the computation graph, producing the
complete gradient $\nabla f \in \mathbb{R}^n$ in a single
backward pass regardless of $n$.

\paragraph{The computational graph.}

Every numerical computation can be decomposed into a sequence
of elementary operations --- additions, multiplications, calls
to \texttt{sin} or \texttt{sqrt} --- each producing an
intermediate variable.  These operations form a directed
acyclic graph (DAG): inputs at the leaves, the final output
at the root, and elementary operations at the internal nodes.
Forward-mode AD traverses this graph from leaves to root,
accumulating derivatives via the chain rule at each node.
The dual-number arithmetic of \cref{sec:ad-forward}
implements exactly this traversal: each intermediate result
carries its derivative forward through the graph.

Reverse-mode AD traverses the same graph in the opposite
direction.  It first performs a forward pass that evaluates
$f$ and records every intermediate variable and the local
partial derivatives at each node.  It then performs a
backward pass from the root to the leaves, accumulating the
chain rule in reverse.  The quantity propagated backward is
the \emph{adjoint} $\bar{v}_i = \partial f / \partial v_i$
--- the sensitivity of the final output to each intermediate
variable $v_i$.

\paragraph{The reverse accumulation rule.}

For a computation that produces intermediate variables
$v_1, \ldots, v_N$ from inputs $x_1, \ldots, x_n$, with the
final output $f = v_N$, the adjoint of each variable is
defined as
%
\begin{equation}
  \bar{v}_i = \frac{\partial f}{\partial v_i} \,.
  \label{eq:adjoint-def}
\end{equation}
%
The backward pass computes these adjoints in reverse
topological order.  If node $v_j$ depends on $v_i$ through
the elementary operation $v_j = \phi(\ldots, v_i, \ldots)$,
then $v_i$'s adjoint accumulates a contribution from $v_j$
via the chain rule.  Summing over all successor nodes gives
the complete adjoint:
%
\begin{equation}
  \bar{v}_i
    = \sum_{j \in \operatorname{succ}(i)}
      \bar{v}_j \,
      \frac{\partial v_j}{\partial v_i} \,.
  \label{eq:adjoint-accumulation}
\end{equation}
%
Each intermediate variable collects sensitivity contributions
from every node that directly uses it.  The factors
$\partial v_j / \partial v_i$ are local partial derivatives
of single elementary operations --- the same derivatives that
forward mode would propagate through the dual part.  The
difference is the direction of accumulation: forward mode
starts at the inputs and pushes derivatives toward the output,
while reverse mode starts at the output and pulls sensitivities
back toward the inputs.
\physnote{The forward and reverse modes correspond to two
different parenthesisations of the same chain-rule product.
Neither changes the mathematical result; the difference is
purely computational --- which order minimises the total
number of multiplications.}

\paragraph{Cost comparison.}

The backward pass visits each node once, performing one
multiplication and one addition per edge in the graph.  The
total cost is proportional to the cost of evaluating $f$
itself, regardless of the number of inputs $n$.  More
precisely, the reverse pass costs at most a small constant
multiple (typically $3$--$5\times$) of the forward evaluation
\cite{Griewank2008}.  Contrast this with forward mode, which
requires $n$ forward passes for the full gradient.  The
crossover is clear:
%
\begin{equation}
  \text{cost}(\nabla f) \;\sim\;
  \begin{cases}
    n \times \text{cost}(f) & \text{forward mode,} \\
    c \times \text{cost}(f) & \text{reverse mode,}
  \end{cases}
  \label{eq:ad-cost-comparison}
\end{equation}
%
where $c \approx 3$--$5$ is independent of $n$.  For the
metric function ($n = 4$), forward mode wins: four passes at
$2\times$ each gives $8\times$, versus a single reverse pass
at $3$--$5\times$ plus the memory overhead of storing the
computation graph.  For a scalar objective function of $100$
parameters, reverse mode is faster by a factor of $\sim\!20$.
\warnnote{The reverse pass requires storing every intermediate
variable from the forward pass, so its memory cost scales
with the depth of the computation graph.  For long
integration chains (thousands of time steps), this can be
prohibitive without checkpointing strategies that trade
recomputation for memory \cite{Griewank2008}.}

\paragraph{Connection to backpropagation.}

Reverse-mode AD is mathematically equivalent to the
backpropagation algorithm used to train neural networks
\cite{Baydin2018}.  A neural network is a composition of
affine maps and elementwise nonlinearities; backpropagation
computes the gradient of a scalar loss function with respect
to all network weights by traversing this composition in
reverse.  The AD perspective reveals that backpropagation is
not specific to neural networks --- it applies to any
differentiable computation, including the evaluation of
spacetime metrics, geodesic integrators, and radiative
transfer solvers.
\histnote{Backpropagation was popularised in machine learning
by Rumelhart, Hinton, and Williams in 1986, but the
underlying reverse-mode AD technique dates to Speelpenning
(1980) and Linnainmaa (1970) \cite{Baydin2018}.}

\paragraph{Potential applications in GR.}

The event-horizon codebase does not use reverse-mode AD ---
forward mode is the natural choice for the metric derivative
computation that dominates the integrator's cost
(\cref{sec:ad-christoffels}).  But reverse-mode AD becomes
relevant whenever a scalar quantity depends on many
parameters:

\begin{itemize}
\item \textbf{Parameter estimation.}  Fitting a black hole
  mass $M$, spin $a$, and observer inclination $\iota$ to
  an observed image requires minimising a scalar residual
  over these parameters.  Reverse-mode AD computes the
  gradient of the residual with respect to all parameters
  in a single backward pass through the ray-tracing
  pipeline.

\item \textbf{Differentiable rendering.}  Modern
  differentiable renderers propagate gradients through the
  entire image-formation process --- from spacetime
  parameters through geodesic integration to pixel
  intensities --- enabling gradient-based optimisation of
  scene parameters.  Reverse-mode AD is the enabling
  technology.

\item \textbf{Waveform modelling.}  In gravitational-wave
  data analysis, matched filtering requires waveform
  templates that depend on source parameters (masses, spins,
  eccentricities).  Reverse-mode AD through a numerical
  relativity solver could provide exact gradients for
  accelerating template-bank searches, though the memory
  cost of differentiating through thousands of time steps
  remains a practical challenge.
\end{itemize}

\noindent
These applications share a common structure: a
high-dimensional parameter space, a scalar objective, and a
computation that is already implemented as a composition of
differentiable operations.  Reverse-mode AD turns any such
computation into a gradient oracle at modest cost ---
precisely the structure that forward mode, with its
per-parameter passes, handles inefficiently.
\xrefnote{The application of AD to numerical general
relativity is surveyed in \cite{Liska2018}; the broader
landscape of AD in scientific computing is reviewed in
\cite{Baydin2018}.}
